% Aula 04 --- a base truncada de splines e o custo de cada nó.
Um \emph{spline} é uma função polinomial por partes, colada suavemente nos
\textbf{nós}. A Aula~04 apresenta a base truncada que os constrói, e este
exercício pergunta por que ela funciona --- e quanto custa cada nó.

Um \emph{spline} de grau $g$ com nós $t_1<\dots<t_J$ é combinação linear de
\[
  1,\ x,\ x^2,\ \dots,\ x^{g},\
  (x-t_1)_+^{g},\ \dots,\ (x-t_J)_+^{g},
  \qquad (u)_+=\max(u,0).
\]

\emph{Sobre as letras:} aqui $g$ é o \textbf{grau} e $J$ o \textbf{número de
nós}. As notas usam $k$ e $p$ para essas duas quantidades; trocamos as letras de
propósito, para não colidir com o $k$ do KNN nem com o $p$ de covariáveis. É o
tipo de colisão que a Aula~04 tem e que convém não carregar para uma prova.
\begin{itens}
  \item \pts{0,6} Quantas funções há na base? Escreva $I$ em função de $g$ e $J$,
        e confira no caso cúbico ($g=3$) com $4$ nós.
  \item \pts{1,0} Mostre que $(x-t_j)_+^{g}$ e suas $g-1$ primeiras derivadas se
        anulam em $x=t_j$. Em seguida, tire a conclusão que interessa: por que
        acrescentar essa função à base \emph{não altera} o ajuste à esquerda de
        $t_j$, e por que a emenda no nó sai suave sem que se precise impor
        restrição nenhuma.
  \item \pts{1,0} Use (b) para justificar a frase "cada nó compra exatamente um
        grau de liberdade". Depois compare com a alternativa ingênua: ajustar um
        polinômio de grau $g$ \textbf{independente} em cada um dos $J+1$ pedaços.
        Quantos parâmetros seriam? Onde foi parar a diferença, e o que ela custa?
  \item \pts{0,9} O ajuste é por mínimos quadrados sobre essa base. Isso faz do
        \emph{spline} um \textbf{ajuste linear} no sentido da Aula~02 --- explique
        por quê. E então responda: se ele é linear, o que exatamente o torna um
        método \emph{não paramétrico}?
\end{itens}

\begin{solucao}
\textbf{(a)} São $g+1$ monômios ($1,x,\dots,x^g$) mais $J$ funções truncadas, uma
por nó:
\[
  I = (g+1) + J .
\]
No caso cúbico com $4$ nós, $I=(3+1)+4=\mathbf{8}$ funções na base --- oito
coeficientes a estimar.

\smallskip
\textbf{(b)} Separe os dois lados do nó.

\emph{À esquerda}, $x<t_j$ dá $(x-t_j)_+=0$: a função é identicamente nula num
intervalo inteiro, e uma função identicamente nula tem todas as derivadas nulas
ali.

\emph{À direita}, $x>t_j$ dá $(x-t_j)_+^g=(x-t_j)^g$, cuja $m$-ésima derivada é
\[
  \frac{g!}{(g-m)!}\,(x-t_j)^{g-m}.
\]
Para $m\le g-1$ o expoente $g-m$ é positivo, e essa expressão tende a $0$ quando
$x\downarrow t_j$.

Os limites laterais coincidem, e valem zero, para a função e suas primeiras $g-1$
derivadas. A $g$-ésima é a primeira a saltar: passa de $0$ para $g!$.

Daí as duas conclusões. Como a função é nula à esquerda, somá-la ao modelo
\textbf{não muda nada antes do nó} --- o coeficiente dela é livre para ajustar o
que vem depois sem estragar o que veio antes. E como ela e suas $g-1$ derivadas se
anulam no nó, a emenda é de classe $C^{g-1}$ automaticamente: a suavidade não
precisa ser imposta, ela está embutida na escolha da base.

\smallskip
\textbf{(c)} Cada nó acrescenta \emph{uma} função à base, logo \emph{um}
coeficiente livre, logo um grau de liberdade. E, por (b), esse grau de liberdade é
\textbf{localizado}: só atua à direita de $t_j$. Você compra flexibilidade onde
quer, e paga por uma unidade de cada vez.

A alternativa ingênua seriam $J+1$ pedaços com $g+1$ coeficientes cada:
\[
  (J+1)(g+1)\ \text{parâmetros.}
\]
No caso cúbico com $4$ nós, $5\times 4=20$ --- contra os $8$ da base truncada.

A diferença são as restrições de continuidade. A base truncada já embute que a
função e suas $g-1$ primeiras derivadas são contínuas em cada nó: $g$ restrições
por nó, $gJ$ no total. E a conta fecha:
\[
  (J+1)(g+1)-gJ = g+1+J = I .
\]
Os $12$ parâmetros que sumiram no caso cúbico são exatamente as $3\times 4$
restrições.

O que a alternativa custa: com mais liberdade, ela ajusta \emph{melhor} o treino
--- e produz saltos nos nós, variância alta e extrapolação ruim. Paga-se
flexibilidade com instabilidade, que é a Aula~01 aparecendo de novo com outra
roupa.

\smallskip
\textbf{(d)} Fixadas as funções da base, o modelo é
$g(x)=\sum_{m=1}^{I}\beta_m f_m(x)$: \textbf{linear nos parâmetros} $\beta_m$,
ainda que profundamente não linear em $x$. Basta montar a matriz de delineamento
com $\mathbb{X}_{im}=f_m(x_i)$ e os $\beta_m$ saem pelas equações normais, como na
Aula~02. Vale inclusive a matriz de projeção $\bm H$ --- e, com ela, o atalho do
LOOCV.

O que torna o método não paramétrico não é a forma do ajuste, então. É o fato de
$I$ --- o número de funções da base, isto é, o número de nós --- \textbf{crescer
com $n$}. Com poucos dados usam-se poucos nós e a curva é rígida; com muitos
dados, mais nós, e a classe de funções representáveis se aproxima de qualquer
função suave.

Essa é a distinção que a Aula~04 faz logo no começo: um método paramétrico fixa a
complexidade \emph{antes} de ver os dados; um não paramétrico deixa que ela cresça
com eles. Não é "ter ou não ter parâmetros" --- é \emph{quando} o número deles
fica decidido.
\end{solucao}
